\documentclass[border=10pt]{standalone}
\usepackage{fontspec}
\setmainfont{Noto Serif}
\newfontfamily\bengalifont{Noto Serif Bengali}[Script=Bengali]
\usepackage[Latin,Bengali]{ucharclasses}
\setTransitionsFor{Bengali}{\begingroup\bengalifont}{\endgroup}
\usepackage{newunicodechar}
\newunicodechar{।}{{\bengalifont\char"0964}}
\newunicodechar{॥}{{\bengalifont\char"0965}}

\usepackage[table]{xcolor}
\usepackage{booktabs,tabularx,array}
\definecolor{bapilrule}{HTML}{1F2937}
\definecolor{bapilstripe}{HTML}{F3F4F6}
\renewcommand{\arraystretch}{1.28}
\setlength{\tabcolsep}{5pt}

\begin{document}
\begin{minipage}{17cm}
{\large\bfseries সারণি ৬।}\ {\large শক্তি, দুর্বলতা, সুযোগ ও হুমকি (SWOT)}\par\vspace{7pt}
\small
\rowcolors{2}{bapilstripe}{white}
\begin{tabularx}{\linewidth}{@{}>{\hsize=0.377\hsize\raggedright\arraybackslash}X>{\hsize=1.623\hsize\raggedright\arraybackslash}X@{}}
\arrayrulecolor{bapilrule}\toprule
\rowcolor{white}\textbf{শ্রেণি} & \textbf{উপাদান} \\
\midrule
\textbf{শক্তি} & বিস্তৃত ডিজিটাল সেবা-ব্যবস্থা (জাতীয় তথ্য বাতায়ন, ই-সেবা, ই-নথি); বৃহৎ পরিসরে ডিজিটাল রূপান্তরের অভিজ্ঞতা; বিপুল, তরুণ ও ডিজিটালভাবে সক্রিয় জনগোষ্ঠী; সক্রিয় বাংলা NLP গবেষক সম্প্রদায়; বিদ্যমান জাতীয় এআই কর্মসূচি \\
\textbf{দুর্বলতা} & বিক্ষিপ্ত ও প্রায়ই স্ক্যানকৃত সরকারি নথি; অসামঞ্জস্যপূর্ণ মেটাডেটা ও লাইসেন্স; দুষ্প্রাপ্য কম্পিউটিং; সীমিত এআই নিশ্চয়তা-দক্ষতা; স্বাধীন মূল্যায়ন কম; প্রকল্প-অর্থায়নের ওপর নির্ভরতা \\
\textbf{সুযোগ} & ২০২৬-এর চতুর্থ ত্রৈমাসিক থেকে জাতীয় এআই নীতি প্রণয়ন; নতুন উপাত্ত-আইন ও জাতীয় ডেটা বিনিময়; স্থানীয়ভাবে চালনাযোগ্য উন্মুক্ত-ওজন মডেল; বৃহৎ ভাষা হিসেবে বাংলার সম্ভাবনা; ডিজিটাল জনসম্পদে উন্নয়ন সহযোগীদের আগ্রহ \\
\textbf{হুমকি} & রাজনৈতিক ও প্রাতিষ্ঠানিক পরিবর্তন; ভুল তথ্য ও কৃত্রিম মিডিয়া; সাইবার আক্রমণ ও বিপর্যয়; বিক্রেতার উপর নির্ভরতা; দৃশ্যমান ব্যর্থতায় আস্থাক্ষয়; মেধা পাচার \\
\bottomrule
\end{tabularx}
\par\vspace{3pt}{\scriptsize\color{gray} উৎস: BAPIL কারিগরি ব্লুপ্রিন্ট, সংস্করণ ১.০ (২০২৬)। CC BY 4.0।}
\end{minipage}
\end{document}
